\begingroup
\begin{flushleft}
\bfseries
\fontfamily{phv}
{\Large
Jersey Recognition and Team Affiliation on Soccer Game State Reconstruction}

\mdseries
\end{flushleft}
\normalsize
\let\clearpage\relax
\chapter{RESUMO}
\endgroup

\noindent O desafio SoccerNet de reconstrução do estado de jogo (\emph{Game State Reconstruction},
GSR) consiste em reconstruir, num mapa bidimensional do campo e a partir de vídeo de câmara única,
a posição, a equipa, a função e o número de camisola de cada jogador. Na solução de referência, dois
componentes destacam-se pelo desempenho reduzido, o reconhecimento do número de camisola e a
atribuição da equipa e do lado do campo. Esta dissertação intervém em quatro dos nove módulos do
\emph{pipeline}: substitui o leitor de números, reescreve os módulos de equipa e de lado e estende
o de reidentificação com um descritor de cor, mantendo a deteção, o seguimento, a calibração e a
classificação de função inalterados, de modo que os ganhos possam ser atribuídos sem ambiguidade às
alterações introduzidas.

O reconhecimento do número foi reconstruído em torno de um modelo multimodal de visão e linguagem,
quantizado para caber na memória gráfica disponível e aplicado a recortes da região dorsal obtidos
por estimativa de pose, com as leituras sucessivas de cada \emph{tracklet} consolidadas por votação.
Na atribuição da equipa, o agrupamento por aparência passou a aferir a sua própria qualidade e a
recorrer à cor do equipamento quando a reidentificação não separa as equipas, e o lado do campo é
decidido por uma cascata de critérios que termina na posição do guarda-redes.

A avaliação abrangeu os $159$ vídeos anotados do SoccerNet-GSR, comparando o sistema proposto com a
solução de referência executada nas mesmas condições e sobre os mesmos \emph{tracklets}. A exatidão
do número de camisola subiu de $15.3\%$ para $37.8\%$, com melhoria em $156$ dos
$159$ vídeos; a da equipa e lado do campo subiu de $83.1\%$ para $91.7\%$, e o número de
vídeos em que o módulo de equipa falha, avaliado por inspeção visual, desceu de $30$
para $14$. O GS-HOTA, métrica oficial do desafio, subiu de $28.95\%$ para
$44.00\%$ no conjunto de teste, com ganhos da mesma ordem na validação e no conjunto de treino,
que não interveio no desenvolvimento; a decomposição da métrica por atributo imputa o ganho
inteiramente aos dois atributos alterados. A comparação revelou ainda um defeito de calibração na
configuração distribuída da solução de referência, corrigido antes da comparação.

As experiências controladas mostram que a margem que resta no número de camisola está limitada pela
informação presente nas imagens. Na atribuição da equipa, a inspeção visual caracterizou as falhas
que subsistem, a maior parte devida ao classificador de função, não alterado, e nove a uma
verificação de equilíbrio em falta na auditoria do agrupamento, cuja correção se identifica.

\hspace{3cm}\\
\indent
\textbf{Palavras-chave: } Reconstrução do Estado de Jogo, Reconhecimento do Número de Camisola,
Atribuição da Equipa, Modelos Multimodais, SoccerNet